\clearpage
\fancyhead[L]{\small GAUSS TO VON LINDENAU.}
\begin{center}
\textbf{GAUSS TO VON LINDENAU.}
\par\vspace{4pt}\hrule\vspace{4pt}
\fontsize{10}{12}\selectfont Monatliche Correspondenz. Vol. XX. Pp. 78, 79. July 1809.
\par\vspace{4pt}\hrule\vspace{4pt}
\end{center}
\begin{flushright} Göttingen, 14 August 1809.
\end{flushright}

\par\indent ---\enspace---\enspace--- What I can send you today are a few observations of the new planets that I have made in the last few days. I was particularly keen to look for Pallas in good time. Last year I could not improve the elements for want of good \emph{complete} observations (I still lack more reliable \emph{declinations}, and I rather fear I will be able to use only my own and the ones made in Milan with the equatorial sector). This time we therefore have to make do with the ephemeris Professor \textsc{Bode} calculated for the 1811 yearbook from the elements determined in 1807. --- Even so, there was no trouble finding the planet that way. On 8 August Professor \textsc{Harding} had searched the area in Pisces where it should be, and on 9 August we immediately recognized it in a star of magnitude 9 that had moved noticeably south. On 10 August the recovery was established with complete certainty. Here are the results of my observations on those days.


\begin{center}\fontsize{9}{11}\selectfont\setlength{\tabcolsep}{5pt}\renewcommand{\arraystretch}{1.15}\begin{tabular}{lrrr} 1809 & \shortstack{Mean time\\in Göttingen.} & \shortstack{Apparent right\\ascension of $\PallasSymbol$} & \shortstack{Apparent north\\declination.}\\\hline Aug. 9 & $11^{\text{u}}\;42^{\text{m}}\;48^{\text{s}}$ & $7^\circ\;51'\;44''$ & $2^\circ\;41'\;6''$\\ 10 & $11\;9\;58$ & $7\;50\;35$ & $2\;32\;37$\\
\end{tabular}\end{center}


\SourceLeaf{320}\par\indent This shows that the ephemeris in the astronomical yearbook now gives the declinations well but makes the right ascensions $20'$ too small.

\par\indent On 10 August I also found Ceres. A few observations, still only partly reduced, show at least that my ephemeris in the May issue of the M. C. gives its position to within a minute. Ceres is of magnitude 8 to 9.

\par\indent On the night of the 12th to the 13th, Professor \textsc{Harding} also thought he had already recognized Vesta as a star of magnitude 9 to 10 very near the position given by the ephemeris calculated by Dr. \textsc{Schumacher}. This still needs to be confirmed, and the first clear night will be used for that purpose.

